\chapter{Conclusão}

Esta dissertação investigou a aplicação de seis Modelos de Linguagem de Grande Porte (LLMs) \textit{open-source} --- Llama, Dolphin, Gemma, Mistral, Alpaca e Qwen --- na conversão automática de normas técnicas brasileiras da AEC para o formato RASE, utilizando \textbf{exclusivamente Engenharia de Prompt}. A tarefa foi decomposta em três níveis sucessivos de normalização (N1, N2 e N3) e avaliada em três experimentos (EN1, EN2 e EN1N2) sobre um \textit{dataset} de 79 normas anotadas, com base em treze métricas: nove de similaridade textual (lexicais, semânticas, de distância e voltadas para geração) e quatro de classificação derivadas da matriz de confusão.

\section{Síntese dos Resultados}

A decomposição em três níveis de normalização viabilizou a conversão de normas AEC em representações RASE sem recurso a \textit{Fine-Tuning} ou \textit{Retrieval-Augmented Generation}. A Engenharia de Prompt, isoladamente, foi suficiente para atingir similaridade semântica próxima a 0{,}90 nos melhores cenários (Dolphin em EN1, com BERTimbau~0{,}891 e Multilingual~0{,}913; Llama em N3 isolado, com SBERT~0{,}947). As métricas baseadas em \textit{embeddings} (SBERT, BERTimbau e Multilingual) e o BERTScore demonstraram-se as mais adequadas para validar transformações normativas, ao tolerarem reformulações que preservam o significado original; o BERTimbau apresentou ganho consistente sobre o SBERT, evidenciando o valor de \textit{embeddings} de domínio jurídico em português.

Globalmente, considerando a média sobre as nove métricas de similaridade nos três experimentos, o \textbf{Dolphin} apresentou a melhor média global (0{,}676), impulsionado pelo domínio absoluto em EN1, sendo também a melhor opção em termos de custo-benefício --- aproximadamente noventa vezes mais rápido que o Llama no pipeline encadeado. O \textbf{Llama} liderou em EN2 e em N3 isolado, mas com latência muito maior. Mistral seguiu como terceiro melhor desempenho, com tempos curtos e qualidade próxima dos líderes. Qwen e Alpaca apresentaram desempenho insatisfatório em todas as métricas. As métricas de classificação confirmaram que os operadores Seleção e Exceção são as classes mais difíceis em N2 para todos os modelos, com F1 próxima de zero, indicando a necessidade de prompts mais especializados para esses operadores.

\section{Contribuições}

Como principais contribuições, esta pesquisa apresenta:

\begin{itemize}
    \item A definição operacional de três níveis de normalização sobre a metodologia RASE: \textbf{N1} (segmentação atômica), \textbf{N2} (identificação dos operadores RASE) e \textbf{N3} (estruturação semântica em JSON), com prompts dedicados a cada nível e a cada operador.
    \item Uma comparação sistemática de seis LLMs \textit{open-source} servidos localmente via Ollama, na tarefa de conversão de normas brasileiras da AEC, sob protocolo experimental idêntico (mesmo \textit{dataset}, mesmos prompts e mesmo hardware).
    \item Um conjunto reproduzível de prompts especializados por nível e por operador, disponível em \texttt{mestrado-rase/prompts/}, que pode ser reutilizado e estendido em pesquisas futuras.
    \item Um \textit{dataset} anotado de 79 normas brasileiras, com referência em todos os níveis (N1, N2 e N3) e por operador RASE, viabilizando avaliações replicáveis.
    \item Uma avaliação multimétrica composta por treze métricas (nove de similaridade textual e quatro de classificação), fornecendo um retrato multidimensional do desempenho de cada modelo, com detalhamento por instância em \texttt{metrics/details/} para apoiar reanálises.
\end{itemize}

\section{Limitações}

A pesquisa apresenta limitações que delimitam o alcance das suas conclusões:

\begin{itemize}
    \item O \textit{dataset} é relativamente pequeno (79 normas) e restrito a normas brasileiras de acessibilidade (NBR~9050 e similares); resultados podem variar em outros domínios normativos.
    \item As métricas de classificação dependem da escolha do critério de acerto: limiar sobre similaridade semântica em N1 e N2 (parcialmente arbitrário) e correspondência exata em N3 (severa demais para campos com variações léxicas legítimas). Critérios mais flexíveis em N3 e múltiplos limiares em N1/N2 ficam como linha de aprofundamento.
    \item A classificação dos operadores Seleção e Exceção em N2 obteve F1 próxima de zero para todos os modelos, indicando uma limitação estrutural da formulação atual do prompt N2 que ainda não foi superada por nenhum dos LLMs avaliados.
    \item Não houve aplicação prática em software BIM (validação fim-a-fim em projetos reais); o ciclo completo de verificação automática de conformidade ainda precisa ser fechado.
    \item Os modelos foram executados em uma única estação de trabalho, sem análise de variância entre execuções, sem variação de temperatura e sem variação de \textit{seed}.
\end{itemize}

\section{Trabalhos Futuros}

Como direções de pesquisa futuras, destacam-se:

\begin{itemize}
    \item \textbf{Aplicação em BIM}: integração dos resultados ao Revit ou Dynamo, fechando o ciclo de verificação automática de conformidade em projetos reais.
    \item \textbf{Prompt N2 dedicado a Seleção e Exceção}: redesenho do prompt N2 ou introdução de uma etapa específica para esses dois operadores, dado que ambos colapsaram em F1 próximo de zero para todos os seis modelos avaliados.
    \item \textbf{\textit{Fine-Tuning} e RAG}: explorar essas técnicas como evoluções naturais da abordagem baseada apenas em prompts, comparando os ganhos relativos em qualidade e custo computacional, em especial para os operadores difíceis identificados acima.
    \item \textbf{Ampliação do \textit{dataset}}: incluir um espectro maior de normas brasileiras (NBRs estruturais, hidráulicas e elétricas) e expandir para normas internacionais.
    \item \textbf{Avaliação humana}: complementar as métricas automáticas com avaliações por especialistas da AEC, capazes de capturar nuances semânticas que os \textit{embeddings} não detectam.
    \item \textbf{Estudo de variância}: realizar múltiplas execuções por modelo, variando configurações de temperatura e de \textit{seed}, para quantificar a estabilidade das saídas.
    \item \textbf{Critérios alternativos de classificação em N3}: substituir a correspondência exata por correspondência normalizada (\textit{lowercase}, remoção de \textit{stopwords}, comparação por sinônimos), revisitando as métricas de Acurácia, Precisão, Revocação e F1 em N3.
    \item \textbf{Otimização do pipeline}: investigar estratégias de mitigação da propagação de erros em EN1N2, como reescrita iterativa, \textit{self-consistency} ou roteamento de operadores difíceis para o modelo mais adequado (\textit{ensemble} entre Dolphin e Llama, conforme o nível).
\end{itemize}

Em síntese, a pesquisa demonstrou que LLMs \textit{open-source}, quando combinados com uma decomposição cuidadosa da tarefa e com prompts especializados, constituem uma alternativa viável e acessível para a automação da conversão de normas técnicas da AEC para representações computáveis, abrindo caminho para a verificação automática de conformidade em projetos baseados em BIM.
